% Propietario conservado: /Users/ruben/Documents/New project/output/EXTREMA_MEDIA_RAZON_RECIPROCIDAD_ES_EN_20260918/ARTICULO/ES/source/sections/orientaciones_reticulares.tex
\subsection{Orientations de support complet et équivariance réticulaire}
\label{x_reciprocidad:paper:24-orientaciones}

L’énumérateur de poids de \eqref{x_reciprocidad:exc:enumerador} contient
\(24y^{12}\). Par conséquent, l’ensemble
\[
 C_W^\times:=
 \{u\in C_W:u_b\in\{1,2\}\text{ pour tout }b\}
\]
possède 24 éléments. Pour chaque \(u\in C_W^\times\), nous définissons
\[
 P_b(u)=
 \begin{cases}
  \mathsf R,&u_b=1,\\
  I,&u_b=2,
 \end{cases}
 \qquad
 g_u:=\bigoplus_{b=0}^{11}P_b(u)\mathsf C P_b(u)^{-1}.
\]

\begin{proposition}[Robustesse des orientations à support total]
\label{x_reciprocidad:paper:robustez-orientaciones}
Pour toute \(u\in C_W^\times\),
\[
 \frac{g_uv-v}{3},
 \quad
 \frac{g_u^{-1}v-v}{3}
 \in N_0.
\]
En particulier, chaque \(g_u\) préserve le voisin construit avec
l’orientation transportée.
\end{proposition}

\begin{proof}
En coordonnées de racines simples,
\[
 \mathsf C\rho-\rho=(-2,-1),
 \qquad
 \mathsf C^{-1}\rho-\rho=(-1,-2).
\]
Les classes discriminantes des deux déplacements sont \(u\) et \(-u\).
Pour le coefficient radial \(a_b\in\{4,1\}\),
\[
 \left\langle
 \frac{(\mathsf C^{\pm1}-I)a_b\rho}{3},a_b\rho
 \right\rangle=-a_b^2.
\]
La somme sur les douze facteurs donne de nouveau \(-27\). La preuve du
\cref{x_reciprocidad:km:estabilidad-vecino} s’applique composante par
composante.
\end{proof}

Pour typer le recalibrage, soit
\(h\in\operatorname{Aut}(C_W)\) monomial : une permutation \(\sigma\) des
douze coordonnées suivie de multiplicateurs
\(\epsilon_b\in\{1,2\}\). Nous définissons son relèvement réticulaire
\begin{equation}
 \widetilde h
 :=P_\sigma\circ
 \bigoplus_{b=0}^{11}\mathsf R^{\nu_b},
 \qquad
 \nu_b=
 \begin{cases}
 0,&\epsilon_b=1,\\
 1,&\epsilon_b=2.
 \end{cases}
 \label{x_reciprocidad:paper:elevacion-monomial}
\end{equation}
Son action sur le discriminant est exactement \(h\). Si le changement transporte
simultanément le code, le drapeau, l’origine incidencielle, le vecteur \(v\), le voisin et les
repères, alors
\begin{equation}
 g_{hu}=\widetilde h\,g_u\,\widetilde h^{-1}.
 \label{x_reciprocidad:paper:naturalidad-isometrias}
\end{equation}
La matrice ternaire \(h\) n’agit sans élévation ni sur \(A_2^{12}\) ni sur
Leech.

